\section{Related Work}

\subsection{Task-specific decoders and EEG foundation models}

Early EEG decoders established convolution as an effective inductive bias for
task-specific analysis. EEGNet combines compact temporal filtering with
depthwise spatial convolution \citep{lawhern2018eegnet}, while DeepConvNet and
ShallowConvNet use stacked temporal and spatial filters with pooling
\citep{schirrmeister2017deep}. Although these models can perform strongly on
their target tasks, they are typically trained from scratch and assume a fixed
channel layout, which limits transfer across datasets, montages, and label
spaces.

EEG foundation models have evolved from contrastive representation learning
toward increasingly large-scale masked, generative, and autoregressive
pretraining. BENDR demonstrated transformer-based representation learning from
raw EEG with a contrastive self-supervised objective
\citep{kostas2021bendr}. BIOT extended large-scale pretraining to heterogeneous
biosignals through frequency-domain tokens \citep{yang2023biot}, while LaBraM
applied masked EEG modeling at scale \citep{jiang2024labram}. Recent
generative and autoregressive efforts, including BrainGPT and Neuro-GPT,
further broadened the design space of EEG foundation models
\citep{yue2024braingpt,cui2024neurogpt}. CBraMod combines time-domain patches,
Fourier magnitude, and criss--cross attention in a compact encoder
\citep{wang2025cbramod}, whereas REVE incorporates physical electrode
coordinates to support large-scale pretraining and unseen montages
\citep{elouahidi2025reve}.

However, EEG foundation models typically require large pretraining corpora and
substantial computational resources. Their architectures and self-supervised
objectives are also often inspired by successful approaches in vision, while
the pretraining data, backbone, tokenization, and input geometry change
simultaneously. This makes it difficult to determine how much of the
downstream gain comes from EEG-specific pretraining itself. Recent work raises
a related concern by showing that carefully tuned conventional decoders remain
competitive across diverse BCI tasks \citep{yang2026are}. We take a
complementary perspective: instead of introducing another large EEG
pretraining pipeline, we retain a mature visual backbone, provide EEG with a
lossless geometry adapter, and test whether visual representations learned
without EEG data can transfer to downstream EEG tasks.

\subsection{Vision transfer and lossless rearrangement}

ImageNet pretraining provides transferable low- and mid-level filters
\citep{deng2009imagenet}. EfficientNet jointly scales width, depth, and
resolution \citep{tan2019efficientnet}, and vision transformers offer a
patch-token alternative \citep{dosovitskiy2021vit}. Generic vision models have
been applied to spectrograms and topographic EEG maps, but those
representations introduce a Fourier transform, interpolation, or an electrode
layout; our input remains the raw waveform and changes indices only.

Pixel shuffle and space-to-depth operations show that a tensor reindexing can
substantially alter the computation of a fixed convolutional architecture
\citep{shi2016realtime}. Temporal folding is related in spirit but asymmetric:
it trades temporal width for a joint electrode--phase height, with each folded
column holding $P$ consecutive samples and each row following one phase at
stride $P$. This explicit signal meaning lets us analyze the original temporal
offsets seen by local 2D neighbors.
